\noindent\textbf{\LARGE{01. Memory}}

\vspace{0.5cm}

\begin{onehalfspace}
\noindent\large\begin{minipage}{\textwidth}

Before anything else, let's establish what memory actually is.

\vspace{0.2cm}

From your program's perspective, memory is just a very long sequence of bytes, each with a unique number called \textbf{address}. On a 64-bit system, addresses are 64-bit numbers, giving a a theoretical address space of $2^{64}$ bytes. On a 32-bit system, addresses are 32-bit numbers.

\vspace{0.2cm}

Think of memory as a very long street of numbered houses. Each house holds exactly one byte. The house number is the address. When your program stores a value, it's putting data into one or more of those houses. When it reads a value, it's looking up a house by number and reading what's inside.

\vspace{0.2cm}

\begin{lstlisting}
int x = 42;
// prints 42
printf("value: %d\n", x);
// prints something like 0x7ffd3a2c1b4c
printf("address: %p\n", (void*)&x);
\end{lstlisting}

\vspace{0.2cm}

The \texttt{\&} operator gives you the address (house number) of a variable. On most systems, an \texttt{int} occupies 4 consecutive bytes, so \texttt{x} occupies 4 houses starting at that address.

\end{minipage}
\end{onehalfspace}
